\chapter{Verificación y Validación: ¿Cómo Confiar en tu Simulación?}
\label{app:vv}

Una simulación cuyos resultados no son fiables es peor que inútil; es peligrosa. Puede llevar a decisiones de diseño incorrectas con consecuencias graves. Por ello, la práctica profesional de la simulación se rige por un marco riguroso conocido como \textbf{Verificación y Validación (V\&V)}. Este apéndice introduce los conceptos fundamentales de este marco.

%==================================================================================
\section{Los Dos Pilares de la Confianza}
%==================================================================================
La V\&V se compone de dos procesos distintos pero complementarios, que responden a dos preguntas diferentes.

%----------------------------------------------------------------------------------
\subsection{Verificación: ¿Resolvemos las Ecuaciones Correctamente?}
%----------------------------------------------------------------------------------
La \textbf{verificación} es un proceso primordialmente \textbf{matemático}. Su objetivo es asegurar que la implementación del modelo (nuestro código FEniCS) resuelve de forma precisa las ecuaciones matemáticas que hemos planteado. No nos preguntamos si la física es correcta, solo si nuestro programa está libre de errores y si la solución numérica es una buena aproximación de la solución matemática exacta (que generalmente no conocemos).

Las actividades de verificación incluyen:
\begin{itemize}
    \item \textbf{Verificación del Código:} Pruebas unitarias, revisión de la implementación, comparación con otros códigos.
    \item \textbf{Verificación de la Solución:} Asegurarse de que la solución numérica converge hacia la solución matemática a medida que refinamos nuestros parámetros numéricos (malla, paso de tiempo).
\end{itemize}

%----------------------------------------------------------------------------------
\subsection{Validación: ¿Resolvemos las Ecuaciones Correctas?}
%----------------------------------------------------------------------------------
La \textbf{validación} es un proceso primordialmente \textbf{físico}. Su objetivo es determinar el grado en que nuestro modelo matemático (las EDPs y sus condiciones de contorno) es una representación precisa de la realidad. En la validación, comparamos los resultados de la simulación con \textbf{datos experimentales}.

Las actividades de validación incluyen:
\begin{itemize}
    \item Diseño de experimentos físicos para medir las cantidades de interés.
    \item Cuantificación de las incertidumbres tanto en las mediciones experimentales como en los parámetros del modelo de simulación.
    \item Comparación estadística entre los resultados de la simulación y los datos del mundo real para determinar si el modelo es ``válido'' para el propósito previsto.
\end{itemize}
La validación responde a la pregunta: ``¿Captura mi modelo la física relevante con la suficiente precisión para ser útil?''.

%==================================================================================
\section{Una Técnica de Verificación Esencial: El Estudio de Convergencia de Malla}
%==================================================================================
La técnica de verificación más fundamental y no negociable en cualquier análisis de elementos finitos es el \textbf{estudio de convergencia de malla} (o refinamiento de malla).

El método de elementos finitos es una aproximación. Su precisión depende directamente del tamaño de los elementos de la malla: cuanto más pequeños sean los elementos (una malla más ``fina''), más precisa debería ser la solución. Sin embargo, una malla más fina también implica un mayor costo computacional.

El objetivo de un estudio de convergencia es encontrar el punto óptimo donde la malla es lo suficientemente fina para que los resultados sean independientes de ella, pero no tan fina como para desperdiciar recursos computacionales.

%----------------------------------------------------------------------------------
\subsection{El Procedimiento}
%----------------------------------------------------------------------------------
El proceso es sistemático:
\begin{enumerate}
    \item Se elige una cantidad de interés (QoI) escalar y local, que sea importante para el diseño. Por ejemplo, la temperatura máxima, el esfuerzo máximo o el desplazamiento en un punto.
    \item Se resuelve el problema en una secuencia de mallas cada vez más finas. Por ejemplo, con resoluciones de \(16\times16\), \(32\times32\), \(64\times64\), y \(128\times128\) elementos.
    \item Se registra el valor de la QoI para cada malla.
    \item Se grafica la QoI en función del tamaño del elemento (\(h\)) o del número de grados de libertad de la malla.
\end{enumerate}

%----------------------------------------------------------------------------------
\subsection{Interpretación de los Resultados}
%----------------------------------------------------------------------------------
El resultado esperado es que, a medida que la malla se refina, el valor de la QoI converja asintóticamente hacia un valor estable.

\begin{figure}[h!]
    \centering
    % Aquí iría un gráfico de ejemplo de convergencia de malla.
    % Por ahora, usamos un marcador de posición.
    \fbox{\begin{minipage}{0.7\textwidth}
    \centering
    \vspace{5cm}
    \textbf{Gráfico de Ejemplo de Convergencia de Malla} \\
    (Eje Y: Temperatura Máxima, Eje X: Número de Elementos)
    \vspace{5cm}
    \end{minipage}}
    \caption{Ejemplo de un estudio de convergencia. A medida que el número de elementos aumenta, la solución para la temperatura máxima se estabiliza, indicando que se ha alcanzado la independencia de la malla.}
    \label{fig:mesh-convergence}
\end{figure}

Si la solución cambia significativamente con cada refinamiento, significa que las mallas iniciales eran demasiado gruesas y los resultados no eran fiables (\textbf{solución dependiente de la malla}). Si la solución se estabiliza, podemos confiar en que hemos eliminado el error de discretización espacial como una fuente significativa de incertidumbre.

Realizar un estudio de convergencia de malla es el primer paso para pasar de ser un usuario de software a un analista de simulación responsable.